% Bilaga E: vidare läsning. Grunderna först, sedan protokollreferensen, standarderna, verktygen och
% fyra riktningar vidare i ämnet.

\chapter{Vidare läsning}
\label{app:reading}

Den här boken slutar med en CAN-nod som en mikrokontroller kan prata med, och den slutar där med
flit: konstruktionen är förenklad på flera medvetna punkter (\secref{c19:sec:gaps}), och drivrutinen
på andra sidan SPI-länken skrivs av någon annan. Den här bilagan säger var grunderna boken
förutsätter kan repeteras, var protokollet självt står beskrivet, vilka standarder konstruktionen
förhåller sig till, och vart man går härnäst.

% ------------------------------------------------------------------------------------------
\section{Grunderna}
\label{app:reading:basics}

Boken förutsätter digitalteknikens grunder och repeterar dem bara där språket kräver det. \debook{}
(\url{https://github.com/qrtech-academy/digital-electronics}) är boken att gå till när de behöver
fräschas upp. Den är på engelska, fritt tillgänglig, och slutar ungefär där den här boken börjar:
med en tillståndsmaskin byggd av grindar och vippor, simulerad i samma CircuitVerse.

Sex kapitel plus bilagor, och varje kapitel motsvarar något den här boken tar för givet:

\begin{booktable}{@{}c>{\raggedright\arraybackslash}p{0.44\linewidth}L@{}}
  \thead{Kap.} & \thead{Ämne} & \thead{Här i boken} \\
  \midrule
  1 & Logiska grindar, och transistorerna och fördröjningen i dem & \secref{c1:sec:gates},
    \secref{c5:sec:fmax} \\
  2 & Boolesk algebra, De Morgan, talsystem och koder & \secref{c1:sec:boolean} \\
  3 & Karnaughdiagram och minimering & \secref{c2:sec:kmap}, \secref{c8:sec:byhand} \\
  4 & Kombinatoriska byggblock: multiplexer, avkodare, komparator, adderare & \secref{c2:sec:mux} \\
  5 & Vippor, register, metastabilitet och synkronisering & \chapref{c3:ch} och \ref{c4:ch} \\
  6 & Räknare, timers och tillståndsmaskiner & \chapref{c6:ch}, \ref{c7:ch} och \ref{c8:ch} \\
  A & Svar på övningarna & \textemdash{} \\
  B & CircuitVerse & \secref{c1:sec:cv} \\
\end{booktable}

Bilagorna efter B är två skrivningar med lösningsförslag. De är det snabbaste sättet att pröva om
förkunskaperna sitter innan \chapref{c1:ch}, och att skriva en av dem på tid visar tydligare än en
genomläsning vilka av kapitlen ovan som är värda en repetition.

% ------------------------------------------------------------------------------------------
\section{Protokollreferensen}
\label{app:reading:protocol}

\canbook{} är den bok den här boken lutar sig mot varje gång CAN som \emph{protokoll} kommer på tal.
Den är kort och skriven för båda halvorna av ett CAN-system: den som bygger kontrollern i hårdvara
och den som skriver drivrutinen i mjukvara. Den nämner därför varken kurs eller klass, och den är
avsiktligt oberoende av vilken av de två kurserna som läser den.

Den finns i två utgåvor med samma kapitel- och avsnittsnumrering, så varje hänvisning i den här
boken gäller båda: på svenska på \url{https://github.com/Yrgo-26/can-book/blob/main/sv/can-sv.pdf},
och på engelska, som \textenglish{\emph{CAN: the bus, the frame and the controller}}, på
\url{https://github.com/Yrgo-26/can-book/blob/main/en/can-en.pdf}.

Sju kapitel plus ett appendix, och kapitelindelningen är den här bokens del~\ref{part:can} sedd från
protokollhållet:

\begin{booktable}{@{}c>{\raggedright\arraybackslash}p{0.44\linewidth}L@{}}
  \thead{Kap.} & \thead{Ämne} & \thead{Här i boken} \\
  \midrule
  1 & Bussen: topologi, dominant och recessiv, wired-AND, terminering & \secref{c10:sec:can} \\
  2 & Framen: fälten, identifieraren, RTR, DLC, kvittensen & \secref{c10:sec:frameformat} \\
  3 & Bitstoppning och CRC-15 & \secref{c10:sec:stuffing}, \chapref{c13:ch}, \ref{c14:ch} och
    \ref{c15:ch} \\
  4 & Arbitrering, bit för bit & \secref{c10:sec:arbitration}, \secref{c17:sec:arbitration} \\
  5 & Bittajming och synkronisering & \secref{c10:sec:bittiming}, \secref{c12:sec:bittimer} \\
  6 & Felhantering: felramar, felräknare, bus-off & \secref{c19:sec:gaps}, som lucka \\
  7 & Den förenklade kontrollern, och de tre begränsningar som når ända upp &
    \secref{c19:sec:gaps}, \secref{app:project:limits}, bilaga~\ref{app:protocol} \\
  A & Svar på samtliga övningar & \textemdash{} \\
\end{booktable}

Tre saker gör den värd att ha uppslagen bredvid den här boken snarare än läst en gång:

\begin{itemize}
  \item \textbf{Kapitel 1--6 är kravspecifikationen, bit för bit,} för den kontroller
    del~\ref{part:can} bygger. Där den här boken säger ”regeln är fem lika bitar i rad”, säger den
    boken varför regeln finns och vad som går sönder utan den.
  \item \textbf{Kapitel 7 beskriver exakt den här konstruktionen}, utifrån: blocken, vad den
    implementerar av standarden, vad den medvetet låter bli, varför ett registerlager behövs, och de
    tre begränsningar som når ända upp i drivrutinens kod. Den som vill ha helheten på fem sidor
    innan \chapref{c11:ch}:s arkitektur läser det kapitlet först.
  \item \textbf{Fyra övningar per kapitel, alla med papper och penna, alla besvarade i appendix A.}
    De prövar protokollförståelsen utan att en rad VHDL är inblandad, vilket gör dem till den bästa
    förberedelsen som finns inför \chapref{c20:ch}:s läsuppgifter.
\end{itemize}

Böckerna delar författare och sätts på samma sätt, så en referens ser likadan ut i båda. De numrerar
däremot sina kapitel oberoende av varandra: ”kapitel~3” utan vidare bestämning betyder alltid ett
kapitel i \emph{den här} boken, och en hänvisning till den andra skrivs alltid ut i sin helhet.

% ------------------------------------------------------------------------------------------
\section{Standarder och specifikationer}
\label{app:reading:specs}

Boken förhåller sig hela tiden till riktig CAN utan att vara den, och \canbook{} är mellansteget:
den beskriver standarden i den detaljnivå en konstruktör behöver. Vill du därefter se vad som
faktiskt står i originalet:

\begin{itemize}
  \item \textbf{Bosch CAN Specification 2.0} (del A och B) är grunddokumentet: ramformaten,
    bitstoppningen, CRC-polynomet och arbitreringen. Del A är standardramen med 11-bitars
    identifierare, alltså den här bokens ram; del B lägger till den utökade 29-bitars formen.
  \item \textbf{ISO 11898-1} är standardiseringen av samma sak, med datalänklagret och det fysiska
    signalgränssnittet. ISO 11898-2 är höghastighetstransceivern, alltså den krets som skulle sitta
    mellan \code{tx_bus}/\code{bus_en} och en verklig differentiell buss.
  \item \textbf{IEEE 1076-1993} är VHDL-93, språkstandarden \code{--std=93} väljer. Den är läsvärd
    som uppslagsverk snarare än pärm till pärm; paketen \code{ieee.std_logic_1164} och
    \code{ieee.numeric_std} är de två delar du använt genom hela boken.
\end{itemize}

De fyra punkter där bokens kontroller medvetet avviker från standarden är samlade i
\secref{c19:sec:gaps}, och \secref{app:project:limits} säger vilka tre av dem som syns ända ut i
registerkartan.

% ------------------------------------------------------------------------------------------
\section{Verktygen}
\label{app:reading:tools}

\begin{itemize}
  \item \textbf{GHDL:s dokumentation} täcker betydligt mer än de tre kommandon bilaga~\ref{app:sim}
    behöver: vågformsutmatning till VCD eller GHW, kodtäckning, och de andra
    språkstandardsflaggorna. \textbf{GTKWave} är den vanliga vågformsvisaren att para ihop den med,
    och första gången du felsöker en tajmingbugg genom att \emph{titta} på signalerna i stället för
    att läsa \code{report}-rader är den värd sin installationstid.
  \item \textbf{Quartus Prime Lites egen dokumentation}, i synnerhet tidsanalysen: bilaga
    \ref{app:quartus} nöjer sig med ett \code{create_clock}, medan en riktig konstruktion också
    behöver ingångs- och utgångskrav och undantag för korsande klockdomäner.
  \item \textbf{DE0-CV:s användarmanual} från korttillverkaren har den fullständiga pinntabellen,
    kringkretsarna och kopplingsschemat. Den är facit när \file{DE0_CV.qsf} inte täcker något du vill
    använda.
  \item \textbf{CircuitVerse} är gratis och körs i webbläsaren. Det är fortfarande det snabbaste
    sättet att pröva en liten kombinatorisk idé, och \chapref{c9:ch} och \ref{c20:ch} tillåter det
    av just det skälet.
\end{itemize}

% ------------------------------------------------------------------------------------------
\section{Vidare i digital konstruktion}
\label{app:reading:design}

Det här är en bok om ett system, inte om hela ämnet. Fyra riktningar som ligger nära och som
material i övrigt inte saknas för:

\begin{itemize}
  \item \textbf{Tidsanalys och klockdomäner på allvar.} \Chapref{c4:ch} ger dubbelvippsynkroniseraren
    och \secref{c5:sec:fmax} en första Fmax. Nästa steg är setup- och holdanalys, FIFO:er mellan
    klockdomäner, och gray-kodade pekare \textemdash{} det vill säga vad man gör när en hel buss, och
    inte en ensam bit, ska passera mellan två klockor.
  \item \textbf{Verifiering bortom en handskriven testbänk:} självkontrollerande modeller,
    slumpstimuli och täckningsmått. VHDL:s OSVVM och UVVM är de två vanliga ramverken, och steget
    därifrån till SystemVerilog och UVM är det industrin tar.
  \item \textbf{Andra gränssnitt, med samma metod.} Boken bygger CAN; UART, SPI, I\textsuperscript{2}C
    och Modbus har alla samma form \textemdash{} en bittimer, ett skiftregister, en sekvenserare
    \textemdash{} och att bygga en till av dem med bokens mönster är den snabbaste vägen att märka
    hur mycket av \chapref{c11:ch}:s uppdelning som var allmängods.
  \item \textbf{Processorkonstruktion.} En liten RISC-kärna är nästa naturliga steg för den som tyckt
    att tillståndsmaskinen i \chapref{c16:ch} var det intressantaste: samma byggstenar, en
    instruktionsavkodare i stället för en ramavkodare, och ett registerfilsgränssnitt i stället för
    en buss.
\end{itemize}

Gemensamt för alla fyra är att de vilar på samma fem mönster \chapref{c20:ch} räknar upp. Det är den
egentliga behållningen av den här boken: inte CAN, utan att du kan bygga nästa protokoll också.
